\documentclass[11pt]{article}

\usepackage[margin=1in]{geometry}
\usepackage{amsmath,amssymb}
\usepackage{booktabs}
\usepackage{graphicx}
\usepackage{hyperref}
\hypersetup{hidelinks}

\title{Supplementary Material: Background-Guided Residual Flow Matching for Multi-Constraint Seismic Velocity Model Building}
\author{Anonymous Authors}
\date{}

\begin{document}
\maketitle

\section{S1. Evaluated BG-RFM Specification}

\subsection{Role-Aware Condition Objective}

For sample $i$, modality $m$, and role $k\in\{\mathrm{str},\mathrm{num}\}$, let $q_{i,m}$ denote the loader-provided quality state and $a_{i,m}$ effective availability. For the evaluated OpenFWI loader, a present non-well modality has $q_{i,m}=1$. For the well modality, $q_{i,m}\in\{0,1\}$ indicates whether sampled well columns exist. The quality state does not use target error. Availability is
\begin{equation}
a_{i,m}
=
\mathbf{1}\!\left[\mathrm{mask}_{i,m}q_{i,m}>0\right].
\end{equation}

Each role-specific reliability score $r_{i,m,k}$ is produced by adaptive average pooling of the modality feature followed by a linear layer and sigmoid. The unnormalized routing weight is
\begin{equation}
\widetilde w_{i,m,k}=a_{i,m}r_{i,m,k},
\end{equation}
and weights are normalized across available modalities with denominator clamped at $10^{-6}$. If no modality is available, the fused weights remain zero. Thus inference uses only observations, masks, loader-provided availability/quality state, and learned reliability; no ground-truth-derived reliability score is used.

The evaluated reliability-prior term is
\begin{equation}
\mathcal L_{\mathrm{rel}}
=
\operatorname{mean}_{a_{i,m}=1}
\left[
(r_{i,m,k}-p_{m,k})^2
\right],
\end{equation}
with fixed $(\mathrm{structural},\mathrm{numerical})$ priors
\begin{center}
\begin{tabular}{lcc}
\toprule
Modality & Structural & Numerical \\
\midrule
Migrated image & 0.85 & 0.45 \\
Horizon & 0.95 & 0.10 \\
RMS velocity & 0.55 & 0.90 \\
Well & 0.25 & 0.95 \\
\bottomrule
\end{tabular}
\end{center}

Training-only anchors are constructed from a one-level Haar decomposition of the target velocity. The numerical anchor is the resized LL component. The structural anchor is the $\ell_2$ fusion of resized LH, HL, and HH components. Haar boundary handling is replicate/edge, resizing is nearest-neighbor, and anchors are detached before condition-loss use.

The subset-Symile terms group samples by identical modality-availability pattern. Within each group, same-record aligned projected modality representations together with the corresponding projected anchor form the positive tuple; negatives are generated by within-group batch derangements, using up to $\mathrm{batch\ size}-1$ negatives. Groups containing fewer than two samples are skipped. Anchor calibration uses $\ell_2$-normalized projected modality representations and normalized projected anchors in a symmetric InfoNCE objective; a single-item group falls back to cosine distance.

Two separation terms are active. Structural/numerical orthogonality is the absolute cosine similarity between adaptively pooled structural and numerical role maps. Unique-role separation averages the absolute cosine similarities between the pooled unique representation and the two pooled role representations. No additional within-role auxiliary loss is active.

The complete evaluated condition objective is
\begin{equation}
\begin{split}
\mathcal L_C={}&
\mathcal L_{\mathrm{subsym}}^{\mathrm{str}}
+\mathcal L_{\mathrm{subsym}}^{\mathrm{num}}
+0.5\,\mathcal L_{\mathrm{anchor}}
+0.05\,\mathcal L_{\mathrm{rel}}\\
&+0.05\,\mathcal L_{\mathrm{str,num}}^{\perp}
+0.05\,\mathcal L_{\mathrm{unique,role}}^{\perp}.
\end{split}
\end{equation}
Pairwise structural/numerical losses and frequency-penalty terms have zero weight in the evaluated configuration and are therefore not part of the active objective. Modality dropout is disabled.

\subsection{Background, Residual, and Joint Objectives}

The shared filter is a $5\times5$ stride-one average pool with padding 2; the framework default includes padded values in the averaging denominator. The high-pass complement is $P_H=I-P_L$.

The direct background objective is
\begin{equation}
\mathcal L_{\mathrm{bg}}
=
\|\hat B-P_L(V)\|_1
+0.5\,\|\hat B-P_L(V)\|_2^2.
\end{equation}
Background high-frequency, multiscale, low-residual, and quality-gate penalties have zero weight.

For the residual branch, $\bar B=\operatorname{sg}[\hat B]$ is detached before residual-target construction, background-context encoding, and residual reconstruction losses. With identity codec,
\begin{align}
R_{\bar B}&=V-\bar B,\\
R_t&=(1-t)\xi+tR_{\bar B},\qquad t\sim\mathcal U(0,1),\\
\widetilde R&=R_t+(1-t)v_\theta(R_t,t,[C_{\mathrm{str}};\psi_{\mathrm{bg}}(\bar B)]).
\end{align}
The active residual objective is
\begin{equation}
\mathcal L_R
=
\mathcal L_{\mathrm{FM}}
+\|\bar B+\widetilde R-V\|_1
+0.2\,\|P_L(\widetilde R)-P_L(V-\bar B)\|_1 .
\end{equation}
The quality/rho term has zero weight. The final joint objective is
\begin{equation}
\mathcal L_{\mathrm{joint}}
=
\mathcal L_R+\mathcal L_{\mathrm{bg}}+0.01\,\mathcal L_C.
\end{equation}
Residual-loss gradients do not update the background branch through $\hat B$; the background branch receives direct supervision through $\mathcal L_{\mathrm{bg}}$. Shared condition-encoder parameters remain jointly optimized through the active objectives.

\subsection{Tensor Alignment and Architecture Interfaces}

\begin{table}[t]
\centering
\caption{Principal tensor interfaces in the evaluated BG-RFM implementation.}
\begin{tabular}{lll}
\toprule
Quantity & Input shape & Aligned/output shape \\
\midrule
RMS & $[B,1,1000,70]$ & $[B,32,70,70]$ \\
PSTM & $[B,1,1000,70]$ & $[B,32,70,70]$ \\
Horizon & $[B,1,70,70]$ & $[B,16,70,70]$ \\
Well & $[B,1,70,70]$ & $[B,16,70,70]$ \\
$C_{\mathrm{bg}}$ & role-fused & $[B,32,70,70]$ \\
$C_{\mathrm{str}}$ & role-fused & $[B,32,70,70]$ \\
Background predictor & $C_{\mathrm{bg}}$ & $[B,1,70,70]$ \\
Background context & $\hat B$ & $[B,32,70,70]$ \\
Residual FiLM U-Net & state + condition & $[B,1,70,70]$ \\
\bottomrule
\end{tabular}
\end{table}

RMS processing pools the temporal dimension and bilinearly aligns the representation to $70\times70$. PSTM uses anisotropic downsampling/projection before alignment. Horizon and well streams are already on the depth grid. Fusion occurs only after spatial alignment. The final model uses an identity codec and a $70\times70$ state grid. The background U-Net has 192 hidden channels, the residual FiLM U-Net has 128 hidden channels, and the condition projectors use 128-dimensional contrastive embeddings.

Training samples continuous Flow-Matching time $t\in[0,1]$. The legacy configuration field named \texttt{residual\_num\_train\_timesteps} is not a discrete training-time schedule for the evaluated \texttt{unet\_fm\_film} backend. Inference integrates the vector field with 50 explicit Euler steps.

\subsection{Staged Training}

\begin{table}[t]
\centering
\small
\caption{Four-stage training schedule used to initialize and optimize the evaluated model.}
\begin{tabular}{p{0.19\textwidth}p{0.08\textwidth}p{0.10\textwidth}p{0.12\textwidth}p{0.36\textwidth}}
\toprule
Stage & Epochs & Batch & AdamW LR & Objective / selection \\
\midrule
Condition pretraining & 100 & 150 & $10^{-4}$ & $\mathcal L_C$; validation retrieval top-5 \\
Background pretraining & 100 & 100 & $10^{-4}$ & $\mathcal L_{\mathrm{bg}}$; validation loss \\
Residual pretraining & 100 & 64 & $5\times10^{-5}$ & FM + reconstruction + residual low-pass; validation loss \\
Final joint optimization & 100 & 64/process & grouped & $\mathcal L_{\mathrm{joint}}$; validation loss \\
\bottomrule
\end{tabular}
\end{table}

Background pretraining is initialized from the selected condition checkpoint. Residual pretraining is initialized from the background checkpoint and keeps the codec frozen. Final joint optimization warm-starts from the condition, background, and residual stage checkpoints, uses four-device DDP, AdamW with grouped learning rates, and keeps the identity codec frozen. The joint learning rates are $10^{-6}$ for the encoder backbone, $3\times10^{-6}$ for encoder heads, contrastive parameters, and condition adapters, $10^{-4}$ for the background estimator, and $5\times10^{-5}$ for the residual vector field; no scheduler and zero configured weight decay are used.

All prerequisite and joint stages invoke the same deterministic dataset-splitting implementation using a 70/20/10 split and split seed 42. Training and validation loaders access only their corresponding partitions, and test metrics are not used for checkpoint selection.

\section{S2. Data Characterization and Effective Rank}

\subsection{Modality-to-Target Kernel Alignment}

Figure~2(a) uses RBF centered-kernel alignment (CKA), implemented as centered normalized HSIC, as a data-level dependence proxy between each input modality and either the low-pass background or structural complement. Within each OpenFWI subset, each representation is flattened; modality representations are divided by the square root of their feature count, and the RBF bandwidth is the median non-zero pairwise Euclidean distance for that representation. The plotted summary uses the median CKA over the eight subsets. A fixed seeded permutation of target rows provides the shuffled control. CKA is descriptive dependence evidence and is not interpreted as mutual information or causal attribution.

\subsection{Participation-Ratio Effective Rank}

For the descriptive decomposition $V=B+S$, with $B=P_L(V)$ and $S=V-B$, participation-ratio effective rank is
\begin{equation}
r_{\mathrm{eff}}=
\frac{\left(\sum_j\lambda_j\right)^2}{\sum_j\lambda_j^2},
\end{equation}
where $\{\lambda_j\}$ are nonzero covariance eigenvalues. The analysis uses 1,000 fixed training records per OpenFWI subset and no learned dimensionality reducer.

\begin{table}[t]
\centering
\caption{Participation-ratio effective rank of the chosen low-pass background and structural complement.}
\begin{tabular}{lrr}
\toprule
Subset & Background $B$ & Structure $S$ \\
\midrule
FlatVel-A & 2.02 & 24.26 \\
FlatVel-B & 5.66 & 39.67 \\
CurveVel-A & 2.59 & 58.52 \\
CurveVel-B & 7.53 & 200.89 \\
FlatFault-A & 3.06 & 37.83 \\
FlatFault-B & 7.54 & 289.69 \\
CurveFault-A & 4.24 & 49.54 \\
CurveFault-B & 10.54 & 407.91 \\
\midrule
Median & 4.95 & 54.03 \\
\bottomrule
\end{tabular}
\end{table}

The structural complement has higher effective rank in every subset. This is a data-level characterization of the selected decomposition; it does not establish lower conditional entropy, lower uncertainty, or guaranteed optimization difficulty.

\section{S3. Sensitivity to Observation Removal}

A fixed trained BG-RFM model is evaluated without retraining after removing individual observations or observation groups. Removal is implemented in two coordinated steps: the corresponding observation tensor is replaced/zeroed, and its availability/quality mask state is updated so the routing mechanism does not treat the removed observation as available. These tests measure dependence/sensitivity under distribution shift from the full-input training contract; they do not isolate the causal contribution of each modality.

The authoritative full-input metrics are reported once in the main common-evaluation table and are not duplicated here.

\begin{table}[t]
\centering
\caption{Sensitivity to observation removal without retraining or test-time adaptation.}
\begin{tabular}{lrrr}
\toprule
Observation subset & MAE & RMSE & SSIM \\
\midrule
w/o well log & 0.0168 & 0.0295 & 0.9899 \\
w/o horizon & 0.1011 & 0.1669 & 0.8202 \\
w/o RMS velocity & 0.3947 & 0.4705 & 0.0753 \\
w/o well log + RMS velocity & 0.4684 & 0.5303 & 0.0192 \\
PSTM only & 0.4706 & 0.5356 & 0.0154 \\
\bottomrule
\end{tabular}
\end{table}

The model is much more sensitive to removal of RMS velocity or horizons than to removal of the sparse well channel. Because the synthetic RMS field is derived directly from the target velocity model, the strong RMS dependence is specific to this informative synthetic observation contract and should not be extrapolated to field-derived RMS products.

\section{S4. Condition and Transport Diagnostics}

Cross-modal retrieval is measured in the role-projected embedding spaces using symmetric Recall@1: within fixed same-subset galleries of complete-modality records, a query from one modality retrieves keys from another modality, the correct match is the same record, and the two retrieval directions are averaged. The diagnostic uses galleries of 128 records, ten galleries per subset, with gallery seed 42.

Anchor similarity is the cosine similarity between an $\ell_2$-normalized role-projected modality representation and the corresponding normalized projected training-time anchor, averaged over available modalities for each record. The deranged control applies a strict one-row cyclic derangement with no fixed points. The reported anchor margin is the per-record matched-minus-deranged difference.

Role specialization uses the fused conditions. With the same $5\times5$ low-pass operator, $\eta_B$ is the low-frequency magnitude ratio of $C_{\mathrm{bg}}$, $\eta_S$ is the high-frequency magnitude ratio of $C_{\mathrm{str}}$, and $\chi_{BS}$ is the absolute cosine similarity between adaptively pooled background and structural condition vectors.

For anchor-margin uncertainty, 95\% bootstrap intervals are obtained from 1000 resamples with replacement of the per-record matched-minus-deranged differences; the displayed interval endpoints are the 2.5th and 97.5th percentiles. The background and structural bootstrap seeds are 17 and 29, respectively.

For the transport diagnostic, the unit-variance Gaussian base contributes one unit of expected squared energy per state dimension. The resulting Gaussian-base-averaged straight-path target-energy ratio is
\begin{equation}
q_T=
\frac{d^{-1}\|V-\hat B\|_2^2+1}
{d^{-1}\|V\|_2^2+1},
\qquad d=HW .
\end{equation}
All 33,600 held-out samples satisfy $q_T<1$ in the evaluated analysis; the mean is 0.772 with bootstrap 95\% confidence interval $[0.771,0.773]$. This interval is obtained by 2000 sample-level bootstrap resamples of the per-record $q_T$ values using seed 2027. This is a diagnostic description of the transport target in normalized coordinates, not a theorem that residual learning must be easier or more accurate.

\begin{figure}[t]
\centering
\includegraphics[width=0.95\textwidth]{../AAAI2027/figures/figure3_condition_learning/appendix_full_condition_diagnostics.png}
\caption{Extended condition-routing diagnostics. Cross-modal retrieval uses symmetric Recall@1 in role-projected spaces; anchor calibration compares matched and strictly deranged anchor similarities; role-specialization statistics describe the fused conditions; routing weights are post-training encoder outputs. These quantities are diagnostic rather than causal evidence.}
\end{figure}

\section{S5. Matched-Input Baseline Adaptations}

All four comparison models receive the same five-channel observation layout:
\begin{enumerate}
    \item post-stack migrated image, bilinearly resized to $1000\times70$;
    \item RMS velocity, bilinearly resized to $1000\times70$;
    \item horizon map, resized by nearest-neighbor interpolation;
    \item masked well-velocity values, bilinearly resized; and
    \item the binary well mask, resized by nearest-neighbor interpolation.
\end{enumerate}
The comparison changes each source method only as needed to accept this common observation tensor and the supervised target used by the reconstruction task. These rows should therefore be interpreted as architecture/backbone adaptations rather than reproductions of the original source-native protocols.

\paragraph{Parameter-count convention.}
The main table reports trainable parameters from the retained evaluated implementation. Training-only modules can differ by method. BG-RFM has 10,630,216 trainable inference parameters; its enclosing training wrapper contains 10,648,427 parameters. VelocityGAN's reported 25.59M training count includes generator and discriminator, although the discriminator is not used for inference. An exact recount tied to the evaluated UPFWI-derived checkpoint is not available, so the main table leaves that parameter entry blank rather than presenting an approximate value as equally authoritative.

\paragraph{InversionNet adaptation.}
The OpenFWI InversionNet encoder--decoder constructor is retained and receives the common five-channel tensor instead of the source-native waveform layout. It is trained by direct $\ell_1$ velocity regression for 100 epochs with batch size 64, AdamW at learning rate $10^{-4}$ and weight decay $10^{-4}$, single-device 32-bit precision, and no learning-rate scheduler. The final checkpoint is used for evaluation. The reported 24.41M parameters count the complete adapted predictor.

\paragraph{VelocityGAN adaptation.}
The OpenFWI InversionNet generator and OpenFWI discriminator are retained. The generator receives the common five-channel tensor. Its objective is $100\,\mathcal L_1+\mathcal L_{\mathrm{adv}}$; the configured MSE term has zero weight. Wasserstein gradient penalty has weight 10, and the discriminator is updated with five critic steps per generator-update period. Generator and discriminator use AdamW at $10^{-4}$ with $(\beta_1,\beta_2)=(0,0.9)$ and weight decay $10^{-4}$. Training lasts 100 epochs with batch size 64, and the final checkpoint is evaluated.

\paragraph{UPFWI-derived supervised multimodal backbone.}
The OpenFWI UPFWI network backbone is retained, but the defining unsupervised wave-equation loop of the original UPFWI method is not used. The backbone receives the common five-channel tensor and is trained as a supervised direct velocity regressor with $\ell_1$ loss for 100 epochs, batch size 64, AdamW at $10^{-4}$, weight decay $10^{-4}$, single-device 32-bit precision, and no scheduler. The final checkpoint is used for evaluation. This is why the main paper labels the row \emph{UPFWI-derived supervised multimodal backbone} rather than using the source-method name alone.

\paragraph{Latent U-Net adaptation.}
The latent U-Net topology is retained within a multimodal translation model consisting of a five-channel measurement encoder, a 128-channel $70\times70$ latent representation, a depth-2 latent U-Net with two repeated blocks per level and skip connections, and a velocity decoder. The complete adapted model is trained with direct $\ell_1$ velocity regression for 100 epochs, batch size 300, and four-device DDP. The minimum-validation-loss checkpoint is used for the reported comparison. The 35.12M parameter count covers the complete adapted measurement-encoder/latent-U-Net/velocity-decoder predictor.

\section{S6. Qualitative Selection and Additional Cases}

The qualitative case-selection rule is performance-dependent and uses no visual inspection. Within each OpenFWI subset, the selection score is the BG-RFM SSIM minus the largest SSIM among the adapted comparison pool consisting of InversionNet, VelocityGAN, the UPFWI-derived supervised multimodal backbone, Auto-Linear, and Latent U-Net. The per-subset panels therefore show maximum-SSIM-margin cases and must not be interpreted as typical or representative samples.

For the composite supplementary panels, embedded legacy labels map to the manuscript terminology as follows:
\begin{itemize}
    \item PD-BG-RFM $\rightarrow$ BG-RFM;
    \item InversionNet $\rightarrow$ InversionNet adaptation;
    \item VelocityGAN $\rightarrow$ VelocityGAN adaptation;
    \item UPFWI $\rightarrow$ UPFWI-derived supervised multimodal backbone;
    \item Latent U-Net (Large) $\rightarrow$ Latent U-Net adaptation;
    \item Auto-Linear denotes an additional adapted baseline that is not included in the main comparison table.
\end{itemize}

\begin{figure}[t]
\centering
\includegraphics[width=0.98\textwidth]{../AAAI2027/figures/table2_qualitative/table2_representative_cases_1.png}
\caption{Performance-selected per-subset maximum-SSIM-margin cases, part I. These are advantage-selected examples rather than representative cases.}
\end{figure}

\begin{figure}[t]
\centering
\includegraphics[width=0.98\textwidth]{../AAAI2027/figures/table2_qualitative/table2_representative_cases_2.png}
\caption{Performance-selected per-subset maximum-SSIM-margin cases, part II.}
\end{figure}

\begin{figure}[t]
\centering
\includegraphics[width=0.98\textwidth]{../AAAI2027/figures/table2_qualitative/table2_failure_cases.png}
\caption{Failure cases selected by the most negative predefined low-frequency, high-frequency, and SSIM margins.}
\end{figure}

\section{S7. Evaluation Details}

The common evaluation uses batch size 64, \texttt{shuffle=False}, and the fixed 33,600-record test membership reconstructed from the global split with seed 42. Record membership does not depend on batch grouping. Test well positions use seed 1234.

For BG-RFM, one stochastic trajectory is evaluated per record and no ensemble averaging is performed. Training uses continuous $t\in[0,1]$, while canonical inference integrates 50 explicit Euler steps. The base sampling seed is 2027. For each evaluation batch, the stochastic seed is generated by SHA-256 hashing the base seed together with the ordered record identities and reducing the hash modulo $2^{31}$. This batch-keyed seed controls only the random draw; it does not alter test membership. Deterministic baselines use their direct prediction without a sampling seed.

The evaluator applies no post-hoc clipping. MAE and RMSE are computed on normalized velocity maps for each record and then averaged. In particular,
\begin{equation}
\mathrm{RMSE}
=
\frac{1}{N}\sum_{i=1}^N
\sqrt{\frac{1}{HW}\|\hat V_i-V_i\|_2^2}.
\end{equation}
SSIM is a global per-image statistic rather than a sliding-window implementation:
\begin{equation}
\mathrm{SSIM}_i=
\frac{(2\mu_x\mu_y+C_{1,i})(2\sigma_{xy}+C_{2,i})}
{(\mu_x^2+\mu_y^2+C_{1,i})(\sigma_x^2+\sigma_y^2+C_{2,i})},
\end{equation}
where population variances/covariance are used,
$L_i=\max(\max V_i-\min V_i,1)$,
$C_{1,i}=(0.01L_i)^2$, and $C_{2,i}=(0.03L_i)^2$.
The reported SSIM is the mean over test records.

Checkpoint selection is separated from test evaluation. BG-RFM uses the selected validation-loss checkpoint from the final joint stage. InversionNet, VelocityGAN, and the UPFWI-derived supervised multimodal backbone use their final 100-epoch checkpoints. Latent U-Net uses the checkpoint with minimum validation loss. The retained formulation-study runs in Table~3 use validation-based model selection. Test metrics are not used to select checkpoints.

Table~3 is a supporting formulation-level comparison. Its retained runs share the dataset family, split contract, normalization, target, metric definitions, optimizer family, and validation-based model selection, while exact solver-step equality is not asserted across every formulation.

\end{document}
